\begin{center}\begin{tikzpicture}[font=\small,>=stealth,line join=round]
\draw[gray!40,rounded corners=3pt] (0,0) rectangle (14.2,8.3);\node at (7.1,7.700000000000001) {分式证明二：把负分数送回同一规则也会还原};
\begin{scope}[shift={(3.725,3.9499999999999997)},x=0.75cm,y=0.3833333333333333cm]\draw[->,gray] (-3.5,0)--(4.5,0) node[right]{$u$};\draw[->,gray] (0,-6)--(0,6) node[above]{$v$};\draw[blue,thick,solid,domain=-3:-1.5,samples=180] plot(\x,{(1-\x)/(1+\x)});\draw[blue,thick,solid,domain=-0.7:4,samples=180] plot(\x,{(1-\x)/(1+\x)});\draw[gray,thick,dashed,domain=-3:4,samples=180] plot(\x,{\x});\fill[orange!85!black] (2,-0.3333333333333333) circle(2.2pt);\fill[orange!85!black] (-0.3333333333333333,2) circle(2.2pt);\draw[orange!85!black,->] (2,-.333333)--(-.333333,-.333333)--(-.333333,2)--(2,2);\end{scope}\node[anchor=west,align=left,text width=5.0cm] at (8.4,4.0) {$\Psi(u)=\frac{1-u}{1+u}$\\[9pt]图中 $2\to-\frac13\to2$\\[9pt]蓝曲线为 $v=\Psi(u)$\\[9pt]虚线为 $v=u$\\[9pt]一般需证 $\Psi(\Psi(u))=u$};
\node[text width=12.6cm,align=center] at (7.1,.65) {图中两次代入均合法；字母通分证明覆盖所有 $u\ne-1$。};
\end{tikzpicture}\end{center}
